% Part: sets-functions-relations
% Chapter: relations
% Section: operations

\documentclass[../../../include/open-logic-section]{subfiles}

\begin{document}

\olfileid[pt]{sfr}{rel}{ops}
\olsection{Operações sobre Relações}

É frequentemente útil modificar ou combinar relações. Na
\olref[sfr][rel][ord]{prop:stricttopartial}, consideramos a \emph{união}
de relações, que é apenas a união de duas relações vistas como
conjuntos de pares. Do mesmo modo, na \olref[sfr][rel][ord]{prop:partialtostrict},
consideramos a diferença relativa de relações. Eis algumas
outras operações que podemos efetuar sobre relações.

\begin{defn}\ollabel{relationoperations}
Sejam $R$ e $S$ relações, e $A$ um conjunto qualquer.

A \emph{inversa} de $R$ é $R^{-1} = \Setabs{\tuple{y, x}}{\tuple{x,
    y} \in R}$.

O \emph{produto relativo} de $R$ e $S$ é $(R \mid S) =
\{\tuple{x, z} : \exists y(Rxy \land Syz)\}$.

A \emph{restrição} de $R$ a $A$ é $\funrestrictionto{R}{A}= R
\cap A^2$.

A \emph{aplicação} de $R$ a $A$ é $\funimage{R}{A} = \{y :
(\exists x \in A)Rxy\}$
\end{defn}

\begin{ex}
Seja $S \subseteq \Int^2$ a relação sucessora sobre~$\Int$, isto é,
$S = \Setabs{\tuple{x, y} \in \Int^2}{x + 1 = y}$, de modo que $Sxy$ sse $x + 1 = y$.

$S^{-1}$ é a relação antecessora sobre $\Int$, isto é,
$\Setabs{\tuple{x,y}\in\Int^2}{x -1 =y}$.

$S\mid S$ é
$ \Setabs{\tuple{x,y}\in\Int^2}{x + 2 =y}$

$\funrestrictionto{S}{\Nat}$ é a relação sucessora sobre~$\Nat$.

$\funimage{S}{\{1,2,3\}}$ é $\{2, 3, 4\}$.
\end{ex}

\begin{defn}[Fecho transitivo]Seja $R \subseteq A^2$ uma relação binária.

O \emph{fecho transitivo} de~$R$ é $R^+ = \bigcup_{0 < n \in
\Nat} R^n$, onde definimos recursivamente $R^1 = R$ e $R^{n+1} = R^n
\mid R$.

O \emph{fecho reflexivo transitivo} de $R$ é $R^* = R^+ \cup
\Id{A}$.
\end{defn}

\begin{ex}
Tome a relação sucessora $S \subseteq \Int^2$. $S^2xy$ sse $x + 2 =
y$, $S^3xy$ sse $x + 3 = y$, etc. Logo $S^+xy$ sse $x + n = y$ para algum
$n \geq 1$. Em outras palavras, $S^+xy$ sse $x < y$, e $S^*xy$ sse $x \le
y$.
\end{ex}

\begin{prob}
Mostre que o fecho transitivo de $R$ é de fato transitivo.
\end{prob}

\end{document}
